\documentclass[11pt]{article}
\usepackage{mathblatt}
% Übersicht nach mathe-nachhilfe gemeinsam.md „Übersicht und Serie“; Einheiten nach pruefung.md § 2 (Dreiecke).
\newcommand{\kennung}{P4Z}
\newcommand{\bereich}[1]{\par\addvspace{14pt}\noindent{\large\bfseries #1}\par\nobreak\vspace{4pt}}
\newcommand{\bl}[1]{\par\noindent #1\par\vspace{3pt}}
\newcommand{\blrand}[1]{\par\noindent{\color{mbgrau}#1}\par\vspace{3pt}}
\newcommand{\blhier}[1]{\par\noindent\llap{$\blacktriangleright$\hspace{0.5em}}\textbf{#1}\par\vspace{3pt}}
\newcommand{\stufe}[1]{\par\noindent\hspace*{1.5em}{\small #1}\par\vspace{2pt}}
\begin{document}
\pfheftstil
\fancyfoot[C]{{\footnotesize\color{mbgrau}\kennung\hfill\thepage}}
\pfheftkopf{Dreiecke}{P10}

\bereich{Vorher}
\bl{Quadrat und Wurzel}

\bereich{Dreiecke}
\blhier{Länge mit Pythagoras}
\stufe{Gleichung aufstellen}
\stufe{Hypotenuse gesucht}
\stufe{Kathete gesucht}
\stufe{Dreieck versteckt}
\bl{Seite oder Winkel mit sin, cos, tan}
\blrand{Seite im allgemeinen Dreieck, Sinussatz}
\bl{Winkel ohne Rechnung bestimmen oder begründen}
\blrand{Symmetrie}

\bereich{Weiter}
\bl{Flächeninhalt und Umfang}
\end{document}
